\section{Prudence: Cardinal Virtue and Its Parts}
\label{sec:prudence}

Prudence (Greek \emph{phronesis}, Latin \emph{prudentia}) is simultaneously an intellectual virtue and the first cardinal family. It perfects practical reason for action, while the other moral virtues make the appetites receptive to its rule. Aquinas organizes its integral conditions, species or scopes, annexed virtues, and opposed failures in ST II--II, qq.~47--56, receiving Aristotle's account in \emph{Nicomachean Ethics} VI.

\cardinalgroup{The principal cardinal and personal prudence}

\virtuedossier
  {Prudence}
  {Principal cardinal and practical-intellectual virtue, also personal prudence or \emph{prudentia simpliciter} when contrasted with its social species; Aristotle, \emph{Nicomachean Ethics} VI.5 and VI.12--13; Augustine, \emph{On the Morals of the Catholic Church} I.15.25; ST I--II, qq.~57--58; II--II, q.~47, a.~11; q.~48, a.~1; qq.~49--56; CCC 1806.}
  {The contingent, practicable human good in the concrete: not simply what can be done, but what ought to be done here and now for the right end.}
  {To take counsel, judge the fitting means, and command their timely execution under right reason. Command is principal because practical truth must issue in action.}
  {\textbf{Imprudence} is the special contrary vice (\oppositioncode{K}). Its principal species corrupt counsel through precipitation, judgment through thoughtlessness, or command through inconstancy and negligence (ST II--II, q.~53, a.~2).}
  {No quantitative excess of prudence exists (\oppositioncode{NA}). Craftiness, worldly prudence, and undue solicitude are counterfeits (\oppositioncode{X}) that retain practical calculation while changing the end or rule.}
  {Prudence is neither caution, calculation, cleverness, conscience, nor technical expertise. \emph{Synderesis} supplies the habitual grasp of first practical principles; conscience applies knowledge to a particular act; prudence judges and commands well.}

Prudential reasoning has three principal movements. \textbf{Counsel} investigates possible means; \textbf{judgment} selects what fits the end and circumstances; \textbf{command} applies the judgment in action. Good intention without adequate counsel can be reckless, endless consultation without judgment can be evasive, and correct judgment without command can remain morally sterile. Prudence is therefore practical truth in agreement with right appetite, not simply accurate prediction.

\cardinalgroup{Integral parts of complete prudential action}

\begin{studybox}{The eight integral conditions}
Aquinas's eight integral parts name what complete prudential action needs (ST II--II, q.~49): \textbf{memory} faithful to experienced particulars; \textbf{understanding} of the present practical situation and first principles; \textbf{docility} receptive to teachable authority and experience; \textbf{shrewdness} (\emph{solertia}) able to discover promptly for oneself; \textbf{reason} comparing means; \textbf{foresight} ordering them to the end; \textbf{circumspection} attending to circumstances; and \textbf{caution} avoiding evils entangled with otherwise good action. They are constitutive conditions of one complete habit here, not eight coordinate cardinal virtues.
\end{studybox}

\cardinalgroup{Subjective species: prudence for self and communities}

Prudence differs by the whole whose good practical reason directs. Prudence without a qualifier---personal prudence or \emph{prudentia simpliciter}---orders one's own life; Aquinas distinguishes it specifically from prudence directed to a multitude (ST II--II, q.~47, a.~11; q.~48, a.~1). He then distinguishes four social species by formally different common-good objects (q.~50). Because they are expressly species rather than mere jobs or applications, each receives a counted subordinate dossier. Their shared act of right practical direction explains their unity without erasing their specific objects.

\begin{threestable}{Subjective species}{Whole directed}{Visible treatment}
Personal prudence & One person's whole life and good action & The principal dossier above; one row, not an omitted fifth species\\
Regnative prudence & The political community under directive governance & Distinct subordinate dossier\\
Political prudence & A free participant's action toward the political common good & Distinct subordinate dossier\\
Domestic prudence & The household as an intermediate community & Distinct subordinate dossier\\
Military prudence & The community's rational defense against attack & Distinct subordinate dossier\\
\end{threestable}

\virtuedossier
  {Regnative prudence}
  {Thomistic subjective species of prudence with Aristotle's legislative or architectonic prudence as antecedent; Aristotle, \emph{Nicomachean Ethics} VI.8; ST II--II, q.~48, a.~1; q.~50, a.~1.}
  {The common good of a complete political community insofar as it must be directed through law and authoritative governance.}
  {To frame and command just common measures, coordinating persons, offices, and subordinate goods toward civic flourishing.}
  {A corresponding species of \textbf{imprudence} opposes this social scope, but Aquinas supplies no separate proper noun for it (\oppositioncode{U}).}
  {\textbf{Tyranny} redirects public power to private advantage and is a counterfeit of governance (\oppositioncode{X}), not excessive regnative prudence.}
  {The historical title does not restrict the virtue to monarchy: Aquinas extends its directive logic to rightful forms of government and identifies legislation as its principal act. It differs from legal justice, which executes and orders acts to the common good.}

\virtuedossier
  {Political prudence}
  {Subjective species of prudence; Aristotle, \emph{Nicomachean Ethics} VI.8; ST II--II, q.~48, a.~1; q.~50, a.~2.}
  {The common good of political life as it specifies a free participant's practical self-direction under legitimate common rule.}
  {To deliberate, judge, and act as a responsible member of the community, translating sound common direction into one's own free action.}
  {A corresponding species of \textbf{imprudence} opposes this social scope, but Aquinas supplies no separate proper noun for it (\oppositioncode{U}).}
  {Factionalism subjects the common good to a part and counterfeits political prudence (\oppositioncode{X}); it is not an excess of the virtue.}
  {Political prudence is neither passive conformity nor expert knowledge of policy. It concerns a citizen's self-government toward the common good and remains bounded by justice and higher law.}

\virtuedossier
  {Domestic or economic prudence}
  {Subjective species of prudence; Aristotle, \emph{Nicomachean Ethics} VI.8; ST II--II, q.~48, a.~1; q.~50, a.~3.}
  {The common good of a household as a community of persons whose life requires ordered relationships, responsibilities, and material means.}
  {To govern household choices and resources so that wealth and work serve the good life of the persons and relationships entrusted to the household.}
  {A corresponding species of \textbf{imprudence} opposes this social scope, but Aquinas supplies no separate proper noun for it (\oppositioncode{U}).}
  {Domination and acquisitiveness corrupt the household's end and counterfeit its direction (\oppositioncode{X}); they are not excess prudence.}
  {Domestic prudence is not household finance alone: riches are instruments, not its final end. It is distinct from both personal prudence and political prudence because the household is a whole between person and polity.}

\virtuedossier
  {Military prudence}
  {Thomistic subjective species of prudence; ST II--II, q.~48, a.~1; q.~50, a.~4. Aristotle's \emph{Nicomachean Ethics} III.8 distinguishes apparent civic or military courage from true fortitude, and VI.8 supplies only the wider political-prudence setting, not this named species.}
  {The common good insofar as a political community must rationally direct its defense against hostile attack.}
  {To counsel, plan, command, and coordinate defensive action under just political ends and proportionate means.}
  {A corresponding species of \textbf{imprudence} opposes this social scope, but Aquinas supplies no separate proper noun for it (\oppositioncode{U}). Mere lack of technical skill is not automatically moral vice.}
  {\textbf{Aggression} abandons the defensive common-good object and counterfeits military direction (\oppositioncode{X}); it is not an excess of the virtue.}
  {Military art governs the skilled use of instruments, fortitude executes steadfastly amid danger, and military prudence directs the whole action toward the common good. The virtue cannot make an unjust war or means prudent.}

\Needspace{24\baselineskip}
\cardinalgroup{Potential parts: counsel and judgment}

\virtuedossier
  {Euboulia (good counsel)}
  {Potential part of prudence; Aristotle, \emph{Nicomachean Ethics} VI.9; ST I--II, q.~57, a.~6; II--II, q.~51, aa.~1--2.}
  {Possible means to a right practical end, considered under the work of inquiry.}
  {To investigate fitting means well: toward the right end, by sound steps, with adequate scope, and in due time.}
  {\textbf{Precipitation} or temerity is the named failure of counsel opposed to euboulia (\oppositioncode{K}); it rushes past steps that right inquiry requires.}
  {No dedicated excess-side vice belongs to euboulia (\oppositioncode{U}); vacillation is a different failure to reach judgment, not perfected counsel carried too far.}
  {Good counsel does not itself judge or command. A clever proposal for an evil end is not euboulia, and accidental success does not retroactively make inquiry sound.}

Speed alone neither proves nor disproves good counsel. Experience can make a well-formed judgment swift; a novel, grave, or irreversible choice can require extended inquiry. The measure is the practical end and the inquiry actually needed, not a temperamentally preferred tempo.

\virtuedossier
  {Synesis (sound judgment)}
  {Potential part of prudence; Aristotle, \emph{Nicomachean Ethics} VI.10; ST I--II, q.~57, a.~6; II--II, q.~51, a.~3.}
  {Ordinary matters of action as judged under the common rules of practical reason.}
  {To judge correctly that a proposed action fits the relevant moral principles and circumstances.}
  {\textbf{Thoughtlessness} or inconsiderateness is the named failure of judgment opposed to synesis (\oppositioncode{K}); it neglects what right judgment must consider.}
  {No dedicated excess-side vice belongs to synesis (\oppositioncode{U}); censoriousness misuses judgment about persons rather than perfecting \emph{synesis} beyond measure.}
  {\emph{Synesis} evaluates ordinary cases after counsel; it is not theoretical science, juridical authority, or the final command belonging to prudence.}

\virtuedossier
  {Gnome (higher judgment)}
  {Potential part of prudence; Aristotle, \emph{Nicomachean Ethics} VI.11; ST I--II, q.~57, a.~6; II--II, q.~51, a.~4.}
  {Exceptional practical cases that must be judged by higher principles when an ordinary rule does not adequately reach the concrete good.}
  {To recognize and judge an exceptional case according to the higher reason presupposed by a sound general rule.}
  {\textbf{Thoughtlessness} is the shared named failure of judgment, but Aquinas supplies no dedicated contrary unique to gnome (\oppositioncode{U}). Rigid adherence to a formulation is an explanatory failure, not a source-fixed noun.}
  {\textbf{Arbitrariness} counterfeits higher judgment by abandoning law and reason (\oppositioncode{X}); it is not excessive gnome.}
  {\emph{Gnome} is not a personal exemption, inspiration without reasons, or equity as a juridical act; it perfects judgment wherever the exceptional case requires higher practical principles.}

\cardinalgroup{Prudence, law, and uncertainty}

Prudence works with moral certitude proportioned to contingent action. It does not wait for the demonstrative necessity proper to mathematics, yet it must not rename impulse as discernment. Relevant facts, competent counsel, duties, foreseeable consequences, and the hierarchy of goods all matter. Some uncertainty can remain after the best available inquiry; a prudent command then acts on the strongest practical reasons while remaining teachable about effects and changed circumstances.

No virtue can make an intrinsically wrongful object prudent. Nor does a good end sanctify dishonest means. Prudence integrates object, intention, circumstances, law, and concrete responsibility; it never serves as an escape hatch from them.
